% Regression document: intentionally forces continuation pages.
\documentclass[degree=master,degreetype=academic]{tongjithesis}
\usepackage{booktabs}
\addbibresource{tests/references.bib}
\tongjimathversion{2}
\newcommand{\auditparagraph}{%
这是用于验证正文固定行距和页边距的排版测试文字。每个自然段首行缩进两个汉字，段前段后均不增加额外间距。This is a layout regression paragraph with Latin text. 本段只检验版面，不属于论文正文。\par}
\makeatletter
\newwrite\layoutout
\immediate\openout\layoutout=\jobname.layout
\newcommand{\layoutmeasure}[2]{\immediate\write\layoutout{#1=\the\dimexpr#2\relax}}
\AtBeginDocument{%
  \layoutmeasure{paperwidth}{\paperwidth}\layoutmeasure{paperheight}{\paperheight}%
  \layoutmeasure{left}{1in+\hoffset+\oddsidemargin}%
  \layoutmeasure{evenleft}{1in+\hoffset+\evensidemargin}%
  \layoutmeasure{textwidth}{\textwidth}\layoutmeasure{textheight}{\textheight}%
  \layoutmeasure{top}{1in+\voffset+\topmargin+\headheight+\headsep}%
  \layoutmeasure{headtop}{1in+\voffset+\topmargin}%
  \layoutmeasure{footerbaselinefrombottom}{\paperheight-(1in+\voffset+\topmargin+\headheight+\headsep+\textheight+\footskip)}%
  \layoutmeasure{footerboxbottomfrombottom}{\paperheight-(1in+\voffset+\topmargin+\headheight+\headsep+\textheight+\footskip+2.1bp)}%
  \layoutmeasure{bodyleading}{\baselineskip}%
}
\makeatother
\begin{document}
\frontmatter
\makeatletter
\begin{cabstract}
\auditparagraph\newpage\auditparagraph
\end{cabstract}
\begin{eabstract}
English abstract continuation test.\par\newpage
English abstract second page.
\end{eabstract}
\ckeywords{测试, 行距, 页边距}
\ekeywords{test, layout, margins}
\tongji@makeabstract
\makeatother
\tableofcontents
\mainmatter
\chapter{正文布局测试}
\section{行距检查}
\auditparagraph\auditparagraph\auditparagraph
\clearpage
\auditparagraph\auditparagraph\auditparagraph
\newcount\sectiontest
\sectiontest=0
\loop\advance\sectiontest by1
  \section{目录续页检查\the\sectiontest}
  \subsection{二级目录缩进检查}
  测试段落。
\ifnum\sectiontest<22\repeat
\backmatter
\begin{reference}
\nocite{*}\printbibliography[heading=none]
\end{reference}
\begin{appendix}
\chapter{编号检查}
\begin{equation}1+1=2\end{equation}
\begin{figure}[h]\centering\rule{10mm}{10mm}\caption{附录图编号}\end{figure}
\begin{table}[h]\centering\caption{附录表编号}\begin{tabular}{cc}甲&乙\end{tabular}\end{table}
\end{appendix}
\begin{acknowledgement}\auditparagraph\newpage\auditparagraph\end{acknowledgement}
\begin{resume}\auditparagraph\newpage\auditparagraph\end{resume}
\end{document}
